\chapter{Conclusión}\label{sec:conclusion}

Este trabajo estudió la concatenación de descriptores como una elección entre representaciones individuales, completas y seleccionadas. La evaluación completa reunió 21 descriptores, cuatro datasets y dos clasificadores bajo un protocolo que separa selección interna y evaluación externa. Las cinco estrategias se compararon mediante macro-F1 y exactitud, con rankings de Friedman y Quade y comparaciones por pares ajustadas mediante Holm.

Los resultados muestran beneficios descriptivos de combinar representaciones y reducciones dimensionales al seleccionar subconjuntos. Ninguna estrategia obtiene el mayor desempeño en todos los dominios. Los contrastes principales entre cuatro datasets no permiten afirmar una superioridad estadística general de GFS, Top-$k$, Completa o Homogénea frente a Individual.

La elección de una representación debe considerar el desempeño externo y la dimensionalidad bajo un protocolo común. Como siguientes pasos se plantea ampliar la diversidad de datasets, estudiar el efecto de la normalización y medir de forma homogénea el costo de extracción y clasificación.
